% Volume VIII: Computational Neuroscience Mathematics
% Continuity note for Chapter 45.
% Previous dependencies: Chapter 43's conductance equations; Chapter 44's spike times, point neurons, adaptation, and rates; Chapters 7 and 21's gradient and optimization language.
% New durable concepts: current-based and conductance-based synapses, synaptic kernels, release probability, short-term facilitation and depression, Hebbian rules, covariance learning, Oja normalization, STDP windows, three-factor rules, and plasticity as local or approximate optimization.
% Forward hooks: Chapter 46 uses synaptic and spike-history filters in neural GLMs; Chapter 47 uses plasticity and coding examples; Chapter 49 reuses synaptic weights in recurrent circuits; Chapter 52 compares biological and artificial learning.
% Notation commitments: presynaptic spike times are t_k; synaptic gating variables are s(t); synaptic weight is w; pre- and postsynaptic activities are x and y; learning rate is eta; STDP timing difference is Delta t = t_post - t_pre.
\section{Chapter 45: Synapses, Plasticity, and Learning Rules}
\label{ch:synapses-plasticity-learning-rules}

Neurons compute in circuits because spikes from one cell change the voltage and future spiking of another. A synapse is therefore both a transmission device and a memory device. It filters presynaptic spikes into postsynaptic currents, changes its short-term efficacy from moment to moment, and can change its long-term weight through activity-dependent plasticity.

\paragraph{Chapter problem.}
How do synapses convert spike times into currents, how do their strengths vary over short and long timescales, and when can a biological learning rule be interpreted as an optimization rule?

\paragraph{Section map.}
Section 45.1 models synaptic transmission with current-based and conductance-based kernels. Section 45.2 studies short-term depression and facilitation. Section 45.3 develops Hebbian learning, normalization, and spike-timing-dependent plasticity. Section 45.4 compares local plasticity rules with gradient-based optimization and credit assignment.

\subsection{45.1 Synaptic transmission}

\paragraph{Guiding question.}
How does a presynaptic spike train become a postsynaptic current in a mathematical neuron model?

\paragraph{Beginner-facing intuition.}
A presynaptic spike is an event. A postsynaptic current is a time course. Synaptic transmission is the filter between them. A spike can open receptor channels quickly, then those channels close over milliseconds. The postsynaptic neuron sees a pulse or conductance transient, not an eternal delta function.

Two modeling choices appear repeatedly. A current-based synapse injects a current waveform independent of postsynaptic voltage. A conductance-based synapse opens a conductance, so the current depends on postsynaptic voltage through a driving force. Conductance-based models are more biophysical; current-based models are simpler and often sufficient near a fixed operating voltage.

\paragraph{Formal definitions and notation.}
Let presynaptic spikes occur at times \(t_1,t_2,\ldots\). A spike train can be written
\[
s_{\mathrm{pre}}(t)=\sum_k\delta(t-t_k).
\]
An exponential synaptic gating variable obeys
\[
\frac{ds}{dt}=-\frac{s}{\tau_s}+\sum_k q_k\delta(t-t_k),
\]
where \(\tau_s>0\) is the synaptic time constant and \(q_k\) is the effective release amplitude for spike \(k\). Between spikes, \(s(t)\) decays exponentially. Interpreting the displayed equation distributionally, integrating across a small interval containing \(t_k\) gives the jump convention
\[
s(t_k^+)=s(t_k^-)+q_k,\qquad
\tau_s\dot s=-s\ \text{between spikes}.
\]
Some authors instead write the impulse inside \(\tau_s\dot s\), which changes the units and jump normalization.

A current-based synapse writes
\[
I_{\mathrm{syn}}(t)=w\,s(t),
\]
where \(w\) has current units after the normalization is chosen. A conductance-based synapse writes
\[
\boxed{I_{\mathrm{syn}}(t)=g_{\mathrm{syn}}\,s(t)(V(t)-E_{\mathrm{syn}}).}
\]
In a membrane equation with outward ionic currents subtracted,
\[
C\dot V=I_{\mathrm{ext}}-g_L(V-E_L)-g_{\mathrm{syn}}s(t)(V-E_{\mathrm{syn}}).
\]

\paragraph{Core mechanism: filtering a spike train.}
With the jump convention, one spike at \(t_k\) adds
\[
s_k(t)=q_k e^{-(t-t_k)/\tau_s}H(t-t_k),
\]
where \(H\) is the step function. By linearity, a spike train produces
\[
\boxed{s(t)=\sum_{k:t_k<t}q_k e^{-(t-t_k)/\tau_s}.}
\]
Thus the synapse computes a causal convolution of the spike train with an exponential kernel. This is the same mathematical mechanism as a first-order low-pass filter.

\paragraph{Concrete example.}
Let \(\tau_s=5\) ms and let a single spike at \(t=0\) produce \(q=1\). Then
\[
s(5\ \mathrm{ms})=e^{-1}\approx0.37,\qquad
s(10\ \mathrm{ms})=e^{-2}\approx0.14.
\]
If a second spike occurs at \(t=10\) ms, then immediately after it
\[
s(10^+)=e^{-2}+1\approx1.14.
\]
For a conductance-based excitatory synapse with \(g_{\mathrm{syn}}=2\) nS, \(s=0.5\), \(V=-60\) mV, and \(E_{\mathrm{syn}}=0\) mV,
\[
I_{\mathrm{syn}}=2\ \mathrm{nS}(0.5)(-60\ \mathrm{mV})=-60\ \mathrm{pA}.
\]
Under the outward-current convention, this inward current depolarizes the postsynaptic cell.

\paragraph{Computational neuroscience example.}
Fast AMPA-like excitation can be modeled with \(\tau_s\) of a few milliseconds, while slower NMDA-like or GABA-B-like components need longer time constants and sometimes voltage dependence. Inhibitory GABA-A synapses often have \(E_{\mathrm{syn}}\) near or below rest. If \(E_{\mathrm{syn}}\) is close to resting voltage, opening inhibitory channels can reduce input resistance without much hyperpolarization; this is shunting inhibition.

\paragraph{AI and machine-learning example.}
Event-based neural networks use synaptic filters to turn discrete spikes into continuous state variables. A filtered spike trace can serve as an eligibility trace for learning, an input feature for a downstream unit, or a differentiable approximation to spike history. In graph neural networks and message-passing models, a synapse-like edge weight determines how one node influences another, but biological synapses add temporal filtering and voltage-dependent driving forces that ordinary message passing often omits.

\paragraph{Common misconceptions and failure modes.}
First, excitatory and inhibitory are not properties of the presynaptic spike alone; they depend on receptor type and reversal potential. Second, a conductance-based synapse changes the membrane time constant as well as the mean voltage drive. Third, a synaptic weight is not always a current amplitude; in conductance models it may be a maximal conductance. Fourth, ignoring synaptic delay and time course can be misleading when spike timing or oscillations matter.

\paragraph{Exercises.}
\begin{enumerate}[label=\textbf{45.1.\arabic*.}, leftmargin=*]
\item \textbf{Mechanical.} Write the exponential synaptic trace produced by two spikes at \(t_1=0\) and \(t_2=8\) ms with amplitudes \(q_1=q_2=1\).
\item \textbf{Proof or derivation.} Solve \(\tau_s\dot s=-s\) between spikes and derive the exponential decay factor.
\item \textbf{Numerical/computational.} With \(\tau_s=10\) ms, compute \(s(20\ \mathrm{ms})\) after one unit-amplitude spike at \(t=0\). Then add a second spike at \(t=20\) ms and compute \(s(20^+)\).
\item \textbf{Interpretation.} Why does shunting inhibition reduce responsiveness even when it barely changes resting voltage?
\item \textbf{Misconception/debugging.} A modeler changes \(E_{\mathrm{syn}}\) from \(0\) mV to \(-80\) mV but keeps calling the synapse excitatory. Explain the problem.
\end{enumerate}

\subsection{45.2 Short-term plasticity}

\paragraph{Guiding question.}
Why can the same presynaptic spike have different postsynaptic effects depending on recent spike history?

\paragraph{Beginner-facing intuition.}
Synapses have fast internal state. Vesicle resources can be depleted by recent release, making later spikes weaker. Release probability can be facilitated by residual calcium, making later spikes stronger. These changes happen over tens to hundreds of milliseconds, so they are plasticity, but not long-term learning.

Short-term plasticity turns a synapse into a dynamic filter. A depressing synapse emphasizes stimulus onset and weakens during sustained activity. A facilitating synapse can act as a detector of bursts or high-frequency presynaptic firing.

\paragraph{Formal definitions and notation.}
A minimal Tsodyks-Markram style model uses a resource variable \(R(t)\in[0,1]\) and a utilization variable \(u(t)\in[0,1]\). The released fraction at spike \(k\) is proportional to
\[
A_k=u_k^+ R_k^-,
\]
where \(R_k^-\) is the available resource just before the spike and \(u_k^+\) is utilization just after facilitation at that spike.

Between spikes,
\[
\frac{du}{dt}=\frac{U-u}{\tau_f},
\qquad
\frac{dR}{dt}=\frac{1-R}{\tau_d},
\]
where \(U\) is baseline utilization, \(\tau_f\) is a facilitation recovery time, and \(\tau_d\) is a depression recovery time. At a presynaptic spike,
\[
u^+=u^-+U(1-u^-),
\qquad
R^+=R^-(1-u^+).
\]
The postsynaptic conductance jump is then proportional to \(A_k=u^+R^-\).

\paragraph{Core mechanism: paired-pulse ratio.}
The paired-pulse ratio compares the response to a second spike with the response to a first spike. In a pure depression simplification, take \(u=U\) constant. Starting from \(R_1^-=1\), the first amplitude is
\[
A_1=U.
\]
After the spike,
\[
R_1^+=1-U.
\]
If the second spike arrives after interval \(\Delta\), recovery gives
\[
R_2^-=1-\bigl(1-R_1^+\bigr)e^{-\Delta/\tau_d}
=1-Ue^{-\Delta/\tau_d}.
\]
The second amplitude is
\[
A_2=U R_2^-,
\]
so
\[
\boxed{\frac{A_2}{A_1}=1-Ue^{-\Delta/\tau_d}.}
\]
This ratio is less than \(1\) when depression dominates.

\paragraph{Concrete example.}
Let \(U=0.5\), \(\tau_d=800\) ms, and \(\Delta=50\) ms. Then
\[
\frac{A_2}{A_1}=1-0.5e^{-50/800}
\approx 1-0.5(0.939)=0.531.
\]
The second response is about \(53\%\) of the first. If instead facilitation dominates, \(u\) may increase enough on the second spike that \(A_2/A_1>1\), especially when baseline \(U\) is small and \(\tau_f\) is long.

\paragraph{Computational neuroscience example.}
Short-term depression at thalamocortical or intracortical synapses can make a circuit sensitive to changes in firing rate, because sustained presynaptic activity depletes resources. Facilitation can make a synapse preferentially transmit bursts, because isolated spikes have low release probability but closely spaced spikes raise \(u\). These dynamics alter circuit computations before any long-term weight change occurs.

\paragraph{AI and machine-learning example.}
Dynamic synapses in spiking neural networks provide state-dependent weights. They can implement temporal filtering, novelty detection, and short-term memory without adding a separate recurrent unit. In machine-learning language, \(R(t)\) and \(u(t)\) are hidden states on the edge, not on the node.

\paragraph{Common misconceptions and failure modes.}
First, short-term plasticity is not necessarily learning in the long-term-memory sense; it can vanish within seconds. Second, depression does not mean the synapse is inhibitory; it means the efficacy of the same sign is reduced. Third, facilitation and depression can coexist, so paired-pulse behavior can depend on interspike interval. Fourth, fitting only a paired-pulse ratio may not identify all model parameters uniquely.

\paragraph{Exercises.}
\begin{enumerate}[label=\textbf{45.2.\arabic*.}, leftmargin=*]
\item \textbf{Mechanical.} Define \(u\), \(R\), \(U\), \(\tau_f\), and \(\tau_d\) in the short-term plasticity model.
\item \textbf{Proof or derivation.} Derive the pure-depression paired-pulse ratio \(A_2/A_1=1-Ue^{-\Delta/\tau_d}\).
\item \textbf{Numerical/computational.} Compute \(A_2/A_1\) for \(U=0.2\), \(\tau_d=500\) ms, and \(\Delta=100\) ms under the pure-depression approximation.
\item \textbf{Interpretation.} Which kind of short-term plasticity would help a postsynaptic neuron detect high-frequency bursts, and why?
\item \textbf{Misconception/debugging.} A student observes a smaller second EPSP and concludes the synapse has become permanently weaker. Explain an alternative short-term explanation.
\end{enumerate}

\subsection{45.3 Hebbian learning and STDP}

\paragraph{Guiding question.}
How can correlations and spike timing change synaptic weights through local activity-dependent rules?

\paragraph{Beginner-facing intuition.}
Hebbian learning starts from the idea that a synapse strengthens when presynaptic and postsynaptic activity occur together. But the raw rule ``increase when both are active'' is unstable by itself: positive feedback can make weights grow without bound. Useful plasticity rules need normalization, competition, saturation, or balancing depression.

Spike-timing-dependent plasticity (STDP) refines correlation into timing. If a presynaptic spike tends to occur shortly before a postsynaptic spike, strengthening can reinforce a causal input. If the order is reversed, weakening can reduce ineffective or anti-causal connections.

\paragraph{Formal definitions and notation.}
Let \(x\) be presynaptic activity, \(y\) postsynaptic activity, and \(w\) a synaptic weight. A basic Hebbian update is
\[
\Delta w=\eta xy,
\]
where \(\eta>0\) is a learning rate. A covariance version subtracts means:
\[
\Delta w=\eta (x-\bar x)(y-\bar y).
\]
For a linear postsynaptic unit \(y=w^\top x\), Oja's rule is
\[
\Delta w=\eta\, y(x-yw).
\]
The second term prevents unbounded growth and implements normalization.

For spike times, define
\[
\Delta t=t_{\mathrm{post}}-t_{\mathrm{pre}}.
\]
A standard pair-based STDP window is
\[
\Delta w=
\begin{cases}
A_+ e^{-\Delta t/\tau_+}, & \Delta t>0,\\
-A_- e^{\Delta t/\tau_-}, & \Delta t<0,
\end{cases}
\]
with \(A_+,A_->0\). Pre-before-post timing strengthens; post-before-pre timing weakens in this convention.

\paragraph{Core mechanism: Oja normalization.}
The raw Hebbian rule for a linear neuron is
\[
\Delta w=\eta yx.
\]
If inputs have nonzero variance along a direction, repeated updates can increase \(w\)'s norm indefinitely. Oja's rule subtracts \(\eta y^2w\):
\[
\Delta w=\eta yx-\eta y^2w.
\]
To see the normalization effect, compute the first-order change in squared norm:
\[
\Delta \|w\|^2
\approx 2w^\top\Delta w
=2\eta y\,w^\top x-2\eta y^2 w^\top w.
\]
Since \(y=w^\top x\),
\[
\Delta \|w\|^2
\approx 2\eta y^2(1-\|w\|^2).
\]
Thus if \(\|w\|<1\), the norm tends to grow, and if \(\|w\|>1\), it tends to shrink. This is a local stabilization mechanism.

\paragraph{Concrete example.}
For STDP, let \(A_+=0.01\), \(\tau_+=20\) ms, \(A_-=0.012\), and \(\tau_-=20\) ms. If a presynaptic spike occurs \(10\) ms before a postsynaptic spike, then \(\Delta t=10\) ms and
\[
\Delta w=0.01e^{-10/20}\approx0.0061.
\]
If the postsynaptic spike occurs \(10\) ms before the presynaptic spike, then \(\Delta t=-10\) ms and
\[
\Delta w=-0.012e^{-10/20}\approx-0.0073.
\]
The same temporal separation produces different sign depending on order.

\paragraph{Computational neuroscience example.}
Hebbian and STDP-like rules are used to model receptive-field formation, sequence learning, and experience-dependent map changes. In a feedforward sensory pathway, inputs that consistently precede postsynaptic firing can be strengthened, making the postsynaptic neuron more selective to those input patterns. Biological STDP depends on receptor subtype, dendritic location, inhibition, neuromodulators, and firing regime, so the simple exponential window is a useful abstraction rather than a universal law.

\paragraph{AI and machine-learning example.}
Hebbian updates are local alternatives to backpropagation. Oja's rule is closely related to extracting a principal component from input data. Contrastive and self-supervised learning also strengthen relationships between co-occurring representations, though modern objectives usually include negative samples, normalization, or global losses that are not available at a single biological synapse.

\paragraph{Common misconceptions and failure modes.}
First, ``cells that fire together wire together'' is incomplete because unnormalized Hebbian learning can diverge. Second, STDP is not only about the nearest pair of spikes in all biological settings; triplets, rates, dendrites, and neuromodulators can matter. Third, local correlation does not solve all credit assignment problems. Fourth, a plasticity rule can reproduce a data trend without proving that the brain literally implements that exact formula.

\paragraph{Exercises.}
\begin{enumerate}[label=\textbf{45.3.\arabic*.}, leftmargin=*]
\item \textbf{Mechanical.} For \(x=3\), \(y=2\), \(\eta=0.01\), compute the raw Hebbian update \(\Delta w\).
\item \textbf{Proof or derivation.} Use the first-order norm calculation to show why Oja's rule tends to stabilize \(\|w\|\) near \(1\).
\item \textbf{Numerical/computational.} Compute the STDP update for \(\Delta t=5\) ms and \(\Delta t=-25\) ms using \(A_+=0.02\), \(A_-=0.015\), and \(\tau_+=\tau_-=20\) ms.
\item \textbf{Interpretation.} Why might pre-before-post strengthening support sequence learning?
\item \textbf{Misconception/debugging.} A simulation using \(\Delta w=\eta xy\) has weights that explode. Name two stabilizing modifications and explain their mathematical role.
\end{enumerate}

\subsection{45.4 Plasticity and optimization}

\paragraph{Guiding question.}
When is a synaptic plasticity rule a gradient method, and when is it only a local heuristic?

\paragraph{Beginner-facing intuition.}
Backpropagation changes a weight using an error signal that says how the weight affected a downstream loss. A biological synapse directly observes local quantities such as presynaptic activity, postsynaptic voltage, and perhaps a neuromodulatory signal. The central mathematical question is credit assignment: how does a synapse know whether changing itself would improve a global objective?

Some plasticity rules are exactly gradients of simple objectives. Others are approximations, local surrogates, or useful adaptive mechanisms without a clearly specified global loss. The difference matters because ``learning'' is not automatically ``gradient descent.''

\paragraph{Formal definitions and notation.}
Consider a scalar linear neuron
\[
y=wx,
\]
with target \(y^*\) and squared loss
\[
L(w)=\frac12(y-y^*)^2.
\]
The gradient is
\[
\frac{dL}{dw}=(y-y^*)x,
\]
so gradient descent gives
\[
\Delta w=-\eta (y-y^*)x=\eta(y^*-y)x.
\]
The update needs presynaptic activity \(x\), postsynaptic output \(y\), and a target-derived error \(y^*-y\). A raw Hebbian rule \(\Delta w=\eta xy\) lacks the target error and therefore optimizes a different implicit objective, if any.

A three-factor plasticity rule has the form
\[
\Delta w=\eta\, e(t)\,M(t),
\]
where \(e(t)\) is a local eligibility trace built from pre-post activity and \(M(t)\) is a modulatory signal such as reward prediction error, surprise, or task feedback.

\paragraph{Core derivation: local versus supervised gradient.}
For the squared loss above,
\[
dL=(y-y^*)\,dy,
\qquad
dy=x\,dw.
\]
Substituting gives
\[
dL=(y-y^*)x\,dw.
\]
By the definition of gradient in one dimension,
\[
\frac{dL}{dw}=(y-y^*)x.
\]
The factor \(x\) is locally presynaptic. The output \(y\) may be locally postsynaptic. The target \(y^*\), or the error \(y-y^*\), is not generally local to an arbitrary synapse in a multilayer biological circuit. This is the mathematical credit-assignment gap.

\paragraph{Concrete example.}
Let \(x=2\), \(w=0.5\), so \(y=1\). Suppose \(y^*=3\) and \(\eta=0.1\). The gradient update is
\[
\Delta w=\eta(y^*-y)x=0.1(2)(2)=0.4.
\]
The new weight is \(0.9\). A Hebbian update would be
\[
\Delta w=\eta xy=0.1(2)(1)=0.2.
\]
Both strengthen the synapse here, but for different reasons. If the target were \(y^*=0\), the gradient update would be \(-0.2\), while the Hebbian update would still be \(+0.2\). This shows why local correlation and task error are not interchangeable.

\paragraph{Computational neuroscience example.}
Reward-modulated STDP uses spike timing to create an eligibility trace and dopamine-like signals to gate weight change. The eligibility trace marks synapses that were recently active; the global modulatory signal says whether recent outcomes were better or worse than expected. This can approximate policy-gradient logic under assumptions, but biological systems include delays, nonlinear dendrites, multiple neuromodulators, and homeostatic constraints.

\paragraph{AI and machine-learning example.}
Backpropagation solves credit assignment in artificial networks by applying the chain rule through a known computational graph. Biologically motivated alternatives include feedback alignment, equilibrium propagation, predictive-coding learning rules, local losses, contrastive Hebbian learning, and perturbation-based estimators. Each proposal specifies a different mathematical route for making a weight update depend on downstream performance.

\paragraph{Common misconceptions and failure modes.}
First, a local rule is not automatically biologically plausible; the required variables and timescales must be available at the synapse. Second, a gradient-looking formula may optimize only a simplified objective under restrictive assumptions. Third, biological plasticity includes homeostasis, structural changes, metaplasticity, and neuromodulation, not only fast weight updates. Fourth, backpropagation is not made local by renaming the error signal; the source and transport of that signal must be explained.

\paragraph{Exercises.}
\begin{enumerate}[label=\textbf{45.4.\arabic*.}, leftmargin=*]
\item \textbf{Mechanical.} For \(y=wx\) and \(L=\frac12(y-y^*)^2\), identify which factors in the gradient \((y-y^*)x\) are local to the synapse and which require additional information.
\item \textbf{Proof or derivation.} Derive the gradient descent update for \(w\) in the scalar linear neuron.
\item \textbf{Numerical/computational.} With \(x=4\), \(w=0.25\), \(y^*=0\), and \(\eta=0.05\), compute both the supervised gradient update and the raw Hebbian update.
\item \textbf{Interpretation.} Explain why a reward-modulated eligibility trace can help with delayed feedback.
\item \textbf{Misconception/debugging.} A paper claims a rule is ``just backprop in the brain'' but gives only \(\Delta w=\eta xy\). What credit-assignment information is missing?
\end{enumerate}

\paragraph{Closing bridge.}
Synapses turn spikes into filtered currents and weights into adaptive circuit parameters. Chapter 46 shifts the viewpoint from mechanistic generation to statistical observation: given spike times recorded from neurons, how do we model them as point processes, fit neural GLMs, and check whether the model explains the data?

\clearpage
